% Standalone section for \input into a report using amsmath, amssymb, amsthm.
% This is a research appendix, not a separately compiled report.
\newtheorem{auxsyntaxthm}{Theorem}[section]
\newtheorem{auxsyntaxlem}[auxsyntaxthm]{Lemma}
\newtheorem{auxsyntaxprop}[auxsyntaxthm]{Proposition}

\section{Two reusable arithmetic lemmas}
\label{sec:reusable-arithmetic-lemmas}

This appendix proves two independent tools.  The first permits a direct
equational proof compiler to use a simple injective polynomial numbering
without checking that every term code is in the numbering's image.  The
second recovers an exact power of two from a pair of Pell ranks under
weaker hypotheses than a radix-dependent formulation would suggest.
Neither result, by itself, constructs a new equation for
$\operatorname{Con}(\mathrm{ZFC})$.  No claim of historical priority is made.

\subsection{Injective syntax and conservative extra constants}

Throughout this subsection, $\mathbb N=\{0,1,2,\ldots\}$.  Fix the
first-order equational signature $\Sigma$ consisting of constants
$\epsilon,\iota$, binary function symbols $\mathsf U,\mathsf C$, and unary
function symbols $\mathsf N,\mathsf V$.  A fixed theory $\mathcal E$ is a
set of equations in finitely many variables; its axioms are understood
through substitution instances.  The argument applies to any fixed such
set, and in particular to a finite set.

We use the ordinary equational calculus: substitution instances of the
axioms, reflexivity, symmetry, transitivity, and congruence for each
function symbol.  A proof of a closed equation can be taken to consist
entirely of ground equations.  Indeed, replace every variable occurring
in a given finite proof by the ground constant $\epsilon$.  This preserves
the closed conclusion, sends axiom instances to axiom instances, and
preserves every inference.  This observation concerns a specified
syntactic calculus and requires no semantic completeness theorem.

Define integer polynomials
\begin{align}
 J(a,b)&=(a+b)^2+a+1,\\
 U(a,b)&=4J(a,b),& C(a,b)&=4J(a,b)+1,\\
 N(a)&=4a+2,& V(a)&=4a+3.
\end{align}
The code of $\epsilon$ is $0$, and the code of $\iota$ is $1$.
The code $\langle t\rangle$ of a ground $\Sigma$-term is obtained by
applying the corresponding polynomial at every constructor.

\begin{auxsyntaxlem}[Disjoint constructors and decreasing decoding]
The pairing $J:\mathbb N^2\to\mathbb N$ is injective and satisfies
$J(a,b)>\max(a,b)$.  The four constructor images are pairwise disjoint
and disjoint from $\{0,1\}$.  Every argument of a constructor has a
strictly smaller code than its value.
\end{auxsyntaxlem}
\begin{proof}
Put $s=a+b$.  For a fixed $s$, the value of $J$ lies in
$[s^2+1,s^2+s+1]$.  These intervals are disjoint because
$(s+1)^2+1>s^2+s+1$.  The interval identifies $s$, after which
$a=J-s^2-1$ and $b=s-a$ identify the arguments.  Also
$J(a,b)\ge s^2+1>s$ for $s\ge1$, and the case $s=0$ is immediate.
The four constructors have distinct residues modulo $4$; the two binary
constructors have values at least $4$ and $5$, respectively.  The
remaining assertions follow directly from their formulas.
\end{proof}

Let $\mathcal J$ be the set of natural numbers outside the four
constructor images and $\{0,1\}$.  Enlarge $\Sigma$ by a distinct constant
$\kappa_m$ for each $m\in\mathcal J$, obtaining $\Sigma^+$.  Strong
recursion defines a decoding map
$\delta:\mathbb N\to\operatorname{Term}(\Sigma^+)$: decode a constructor
value using its unique arguments, decode $0,1$ as $\epsilon,\iota$, and
decode $m\in\mathcal J$ as $\kappa_m$.  The preceding lemma proves that
this definition terminates.  It also proves, for example,
\[
 \delta(U(a,b))=\mathsf U(\delta(a),\delta(b)),
 \qquad \delta(N(a))=\mathsf N(\delta(a)),
\]
with analogous identities for the other constructors.  Induction gives
$\delta(\langle t\rangle)=t$ for every original ground term $t$.

For an axiom term $t(x_1,\ldots,x_h)$, let
$\widehat t(n_1,\ldots,n_h)$ be its integer-polynomial evaluation under
these arithmetic constructors.  Arbitrary natural values $n_i$ are
allowed.  Constructor compatibility gives
\[
 \delta\bigl(\widehat t(\boldsymbol n)\bigr)
   =t\bigl(\delta(n_1),\ldots,\delta(n_h)\bigr).
\]
Thus no test that $n_i$ belongs to the image of the original term
numbering is needed.

\begin{auxsyntaxthm}[Conservative interpretation of unused codes]
\label{thm:syntax-retraction}
Let $s=t$ be a closed $\Sigma$-equation.  There exists a finite arithmetic
equational proof of
$(\langle s\rangle,\langle t\rangle)$, using unrestricted natural term
codes, polynomial evaluations of $\mathcal E$-instances, and the encoded
equational rules, if and only if $\mathcal E\vdash s=t$ in the specified
ordinary equational calculus.
\end{auxsyntaxthm}
\begin{proof}
An ordinary ground proof can be encoded line by line, proving the
reverse implication.  Conversely, apply $\delta$ to the two term codes
on every line of an arithmetic proof.  The constructor identities show
that it becomes an ordinary ground proof in the signature $\Sigma^+$.
Every arithmetic axiom instance becomes an actual $\mathcal E$-instance
with ground $\Sigma^+$-terms as substitutions.

Now apply the term homomorphism
$\rho:\operatorname{Term}(\Sigma^+)\to\operatorname{Term}(\Sigma)$
which fixes every original symbol and sends every $\kappa_m$ to
$\epsilon$.  It sends an axiom instance to an axiom instance and
preserves each equational rule.  It fixes the closed terms $s,t$.
Consequently the transformed finite proof is an original-signature
proof of $s=t$.  This is a syntactic retraction argument; it assumes
neither consistency nor arithmetic soundness of any theory.
\end{proof}

\subsubsection*{One explicit polynomial for local inference}

Represent a conclusion by its pair of term codes $(a,b)$ and two premises
by $(c,d)$ and $(e,f)$.  For one additional natural number $z$, put
\begin{align}
 R_0&=(a-b)^2,\\
 T_0&=(c-e)^2+(a-d)^2+(b-f)^2,\\
 B_0&=(a-U(c,e)-z)^2+(b-U(d,f)-z)^2,\\
 W_0&=(a-4c-2-z)^2+(b-4d-2-z)^2,\\
 I(a,b,c,d,e,f,z)&=z^2(z-1)^2+R_0T_0B_0W_0.
 \label{eq:local-inference-poly}
\end{align}
Here $U$ in the two displayed occurrences is the already defined
integer polynomial; these are abbreviations, not extra unknowns.

\begin{auxsyntaxprop}[Exact local rule check]
For natural term codes $a,b,c,d,e,f$, the equation $I=0$ holds precisely
when $z\in\{0,1\}$ and at least one of the following holds:
\begin{enumerate}
 \item $a=b$, the reflexivity rule;
 \item $c=e$, $a=d$, $b=f$, the rule taking $c=d$ and $c=f$ to $d=f$;
 \item $(a,b)=(U(c,e),U(d,f))$ if $z=0$, or
       $(a,b)=(C(c,e),C(d,f))$ if $z=1$;
 \item $(a,b)=(N(c),N(d))$ if $z=0$, or
       $(a,b)=(V(c),V(d))$ if $z=1$.
\end{enumerate}
The total degree of $I$ is exactly $10$.  The variable $z$ is its only
local auxiliary coordinate beyond the six equation ports.
\end{auxsyntaxprop}
\begin{proof}
Every displayed factor and summand is nonnegative.  Thus $I=0$ forces
both $z^2(z-1)^2=0$ and $R_0T_0B_0W_0=0$.  The former is exactly the
conjunction $z\in\{0,1\}$; the latter is exactly the disjunction of
the four rule checks.  Each sum of squares vanishes precisely when all
its summands vanish.  This also proves the converse.

The four factors have degrees $2,2,4,2$, so $\deg I\le10$.
Specializing gives the polynomial identity
\[
 I(t,0,t,0,0,0,0)
 =2t^4\bigl((4t^2+3t+4)^2+16\bigr)
       \bigl((3t+2)^2+4\bigr).
\]
Its leading coefficient is $288$, proving $\deg I=10$.
\end{proof}

The second rule and reflexivity derive symmetry: use the premises
$c=d$ and $c=c$.  They then derive ordinary transitivity by first
reversing $a=b$ and using $b=a$, $b=c$.  Binary congruence of both
arguments includes one-argument congruence by using a reflexive second
premise.  Hence these local rules generate the ordinary ground
equational calculus when supplemented with the axiom instances.

To make the finite-proof convention exact, seed the list by the valid
reflexive equation $(0,0)$.  Each later line is either an axiom instance
or has two earlier premise indices and a natural $z$ satisfying
\eqref{eq:local-inference-poly}.  A one-premise rule may use the seed
as its unused second premise.  This convention excludes circular
justifications, including a purported first inference from itself.
The premise-index restrictions and the unbounded list are part of the
finite-certificate definition; they have \emph{not} been eliminated
into the local polynomial $I$.

\paragraph{What remains to obtain a consistency equation.}
One must supply an explicit finite equational theory, prove the required
bidirectional syntactic consistency translation, materialize its actual
axiom terms, and compile the entire unbounded finite-proof certificate
into a fixed Diophantine equation.  No variable count, degree bound, or
arithmetic-operation improvement for that unconstructed endpoint follows
from the local degree-$10$ calculation.

\subsection{Exact Pell power recovery without an early radix bound}

For an integer $b\ge2$, define the Pell sequences by
\[
 \chi_b(0)=1,\quad \psi_b(0)=0,\qquad
 \begin{pmatrix}\chi_b(m+1)\\\psi_b(m+1)\end{pmatrix}
 =
 \begin{pmatrix}b&b^2-1\\1&b\end{pmatrix}
 \begin{pmatrix}\chi_b(m)\\\psi_b(m)\end{pmatrix}.
\]
In particular, $\psi_b(1)=1$ and
$\psi_b(m+2)=2b\psi_b(m+1)-\psi_b(m)$.

\begin{auxsyntaxlem}[Recurrence bounds]
For every $b\ge2$ and integer $m\ge1$,
\[
 (2b-1)^{m-1}\le\psi_b(m)\le(2b)^{m-1}.
\]
\end{auxsyntaxlem}
\begin{proof}
The sequence $\psi_b(m)$ is nonnegative and strictly increasing from
$m=0$ onward.  This follows inductively from
$\psi_b(0)=0$, $\psi_b(1)=1$ and the second-order recurrence.  For
$m\ge1$, the same recurrence therefore gives
\[
 (2b-1)\psi_b(m)
 \le\psi_b(m+1)
 \le2b\psi_b(m).
\]
Iterating from $\psi_b(1)=1$ proves both bounds.
\end{proof}

\begin{auxsyntaxlem}[Exact quotient identity]
\label{lem:pell-quotient-exact}
Let $a\ge1$, $A=a+2$, and $H=4a+3$.  For every $m\ge0$,
\[
 \chi_A(m)-a\psi_A(m)
 =2^m+H\sum_{j=0}^{m-1}2^{m-1-j}\psi_A(j),
\]
where the empty sum is zero.  In particular,
$\chi_A(m)-a\psi_A(m)\equiv2^m\pmod H$.
\end{auxsyntaxlem}
\begin{proof}
Put $W_m=\chi_A(m)-a\psi_A(m)$ and note that
$A^2-1=a^2+H$.  The defining recurrences give
\[
 W_{m+1}
 =(A-a)\chi_A(m)+(A^2-1-aA)\psi_A(m)
 =2W_m+H\psi_A(m).
\]
Since $W_0=1$, iterating this identity yields the asserted formula.
The sign is that of the positive sequence $\chi_A$; it is not a
statement obtained by replacing a possibly signed Pell coordinate by
its absolute value without further justification.
\end{proof}

\begin{auxsyntaxthm}[Power recovery]
\label{thm:pell-power-recovery-clean}
Let $X,Y$ be positive integers and $n\ge2$.  Put
\[
 R=2n-1,\quad a=(X+1)Y,\quad A=a+2,\quad
 P=2XY^2+1,\quad H=4a+3.
\]
Suppose
\[
 c=\psi_A(R),\qquad k=2\psi_P(n),\qquad
 kY<c<k(Y+1),
\]
and suppose the \emph{positive} main root
$D=\chi_A(R)$ satisfies
\[
 D-ac\equiv X\pmod H.
\]
Then $X=2^R$.  No hypothesis involving a radix $q$, divisibility of $X$
or $Y$ by $q$, or inequalities $X\ge q$, $Y\ge q^3$ is required.
\end{auxsyntaxthm}
\begin{proof}
Both $A$ and $P$ are at least $2$.  Direct calculation gives
\[
 (2A-1)^2-8P
 =4(X-1)^2Y^2+12(X+1)Y+1>0.
\]
Since $R-1=2(n-1)$, the recurrence bounds imply
\begin{align*}
 \frac ck
 &\ge \frac{(2A-1)^{2(n-1)}}{2(2P)^{n-1}}\\
 &=\frac12\left(\frac{(2A-1)^2}{2P}\right)^{n-1}
 >\frac{4^{n-1}}2=2^{R-2}.
\end{align*}
The strict ratio bound $c/k<Y+1$ therefore gives
$Y+1>2^{R-2}$.  Since $R\ge3$ and $Y$ is an integer,
$Y\ge2^{R-2}$.  Consequently
\[
 H=4(X+1)Y+3\ge8Y+3\ge2^{R+1}+3>2^R.
\]
Also $0<X<H$ follows directly from $X,Y\ge1$.

Lemma~\ref{lem:pell-quotient-exact} gives
$D-ac\equiv2^R\pmod H$.  The hypothesis gives the same residue as $X$.
Both $X$ and $2^R$ lie strictly between $0$ and $H$, so they are equal.
\end{proof}

The lower ratio bound $kY<c$ is natural in the intended applications,
although this particular proof uses only the upper bound.  If the source
instead supplies $D=X+ac+H\gamma$, then after power recovery the exact
quotient formula identifies
\[
 \gamma=\sum_{j=0}^{R-1}2^{R-1-j}\psi_A(j)>0,
\]
because $R\ge3$.  Thus the same elementary identity also verifies the
sign of the recovered quotient.  The hypotheses giving the ranks
$R=2n-1$, the positive main root, and the strict ratio remain essential
inputs to this lemma; they are not established by the lemma itself.

\paragraph{Verification boundary.}
The accompanying standard-library script checks constructor decoding and
retraction, local-rule examples and negative controls, the exact
degree-$10$ specialization, recurrence bounds, and the exact Pell
quotient identity.  Those finite checks are reproducible corroboration.
The unrestricted mathematical statements are established by the proofs
above; no finite check establishes a consistency equivalence.
